\documentclass[11pt]{article}
\newcommand{\ListaTitulo}{Lista 14 --- Fatoriais $3^k$, fracionamento e superfície de resposta}
% Preâmbulo compartilhado — MATD48, Listas de Exercícios 2026
% Usado por Lista{NN}.tex e Gabarito{NN}.tex via % Preâmbulo compartilhado — MATD48, Listas de Exercícios 2026
% Usado por Lista{NN}.tex e Gabarito{NN}.tex via % Preâmbulo compartilhado — MATD48, Listas de Exercícios 2026
% Usado por Lista{NN}.tex e Gabarito{NN}.tex via \input{preamble}

\usepackage[utf8]{inputenc}
\usepackage[T1]{fontenc}
\usepackage[brazil]{babel}
\usepackage{fullpage}
\usepackage{amsmath,amssymb}
\usepackage{enumitem}
\usepackage{xcolor}
\usepackage{fancyhdr}
\usepackage{titlesec}
\usepackage{listings}
\usepackage{hyperref}
\usepackage{booktabs}
\usepackage{graphicx}

\definecolor{matd48azul}{HTML}{0B4F6C}
\definecolor{matd48verde}{HTML}{2E8B57}

\titleformat{\section}{\normalfont\Large\bfseries\color{matd48azul}}{\thesection}{1em}{}
\titleformat{\subsection}{\normalfont\large\bfseries\color{matd48azul}}{\thesubsection}{1em}{}

\providecommand{\ListaTitulo}{Lista de Exercícios}

\setlength{\headheight}{14pt}
\pagestyle{fancy}
\fancyhf{}
\lhead{\small MATD48 --- Planejamento de Experimentos A}
\rhead{\small \ListaTitulo}
\cfoot{\thepage}
\renewcommand{\headrulewidth}{0.4pt}

\lstset{
  basicstyle=\ttfamily\small,
  breaklines=true,
  frame=single,
  columns=fullflexible,
  backgroundcolor=\color{gray!8},
  commentstyle=\color{matd48verde},
  keywordstyle=\color{matd48azul}\bfseries,
  language=R,
  extendedchars=true,
  literate=%
    {á}{{\'a}}1 {à}{{\`a}}1 {â}{{\^a}}1 {ã}{{\~a}}1
    {é}{{\'e}}1 {ê}{{\^e}}1 {í}{{\'i}}1 {ó}{{\'o}}1
    {ô}{{\^o}}1 {õ}{{\~o}}1 {ú}{{\'u}}1 {ç}{{\c c}}1
    {Á}{{\'A}}1 {À}{{\`A}}1 {Â}{{\^A}}1 {Ã}{{\~A}}1
    {É}{{\'E}}1 {Ê}{{\^E}}1 {Í}{{\'I}}1 {Ó}{{\'O}}1
    {Ô}{{\^O}}1 {Õ}{{\~O}}1 {Ú}{{\'U}}1 {Ç}{{\c C}}1
}

\hypersetup{colorlinks=true, linkcolor=matd48azul, urlcolor=matd48azul, citecolor=matd48azul}

\newcommand{\cabecalho}[2]{%
  \begin{center}
    {\Large\bfseries\color{matd48azul} #1}\\[0.3em]
    {\large #2}\\[0.3em]
    \rule{\linewidth}{0.6pt}
  \end{center}
  \vspace{0.5em}
}

\newenvironment{questao}[1]{\vspace{1em}\noindent\textbf{Questão #1.}\hspace{0.4em}}{\par}
\newenvironment{solucao}{\vspace{0.3em}\noindent\textit{Solução.}\hspace{0.4em}\color{black!85}}{\par\vspace{0.5em}}


\usepackage[utf8]{inputenc}
\usepackage[T1]{fontenc}
\usepackage[brazil]{babel}
\usepackage{fullpage}
\usepackage{amsmath,amssymb}
\usepackage{enumitem}
\usepackage{xcolor}
\usepackage{fancyhdr}
\usepackage{titlesec}
\usepackage{listings}
\usepackage{hyperref}
\usepackage{booktabs}
\usepackage{graphicx}

\definecolor{matd48azul}{HTML}{0B4F6C}
\definecolor{matd48verde}{HTML}{2E8B57}

\titleformat{\section}{\normalfont\Large\bfseries\color{matd48azul}}{\thesection}{1em}{}
\titleformat{\subsection}{\normalfont\large\bfseries\color{matd48azul}}{\thesubsection}{1em}{}

\providecommand{\ListaTitulo}{Lista de Exercícios}

\setlength{\headheight}{14pt}
\pagestyle{fancy}
\fancyhf{}
\lhead{\small MATD48 --- Planejamento de Experimentos A}
\rhead{\small \ListaTitulo}
\cfoot{\thepage}
\renewcommand{\headrulewidth}{0.4pt}

\lstset{
  basicstyle=\ttfamily\small,
  breaklines=true,
  frame=single,
  columns=fullflexible,
  backgroundcolor=\color{gray!8},
  commentstyle=\color{matd48verde},
  keywordstyle=\color{matd48azul}\bfseries,
  language=R,
  extendedchars=true,
  literate=%
    {á}{{\'a}}1 {à}{{\`a}}1 {â}{{\^a}}1 {ã}{{\~a}}1
    {é}{{\'e}}1 {ê}{{\^e}}1 {í}{{\'i}}1 {ó}{{\'o}}1
    {ô}{{\^o}}1 {õ}{{\~o}}1 {ú}{{\'u}}1 {ç}{{\c c}}1
    {Á}{{\'A}}1 {À}{{\`A}}1 {Â}{{\^A}}1 {Ã}{{\~A}}1
    {É}{{\'E}}1 {Ê}{{\^E}}1 {Í}{{\'I}}1 {Ó}{{\'O}}1
    {Ô}{{\^O}}1 {Õ}{{\~O}}1 {Ú}{{\'U}}1 {Ç}{{\c C}}1
}

\hypersetup{colorlinks=true, linkcolor=matd48azul, urlcolor=matd48azul, citecolor=matd48azul}

\newcommand{\cabecalho}[2]{%
  \begin{center}
    {\Large\bfseries\color{matd48azul} #1}\\[0.3em]
    {\large #2}\\[0.3em]
    \rule{\linewidth}{0.6pt}
  \end{center}
  \vspace{0.5em}
}

\newenvironment{questao}[1]{\vspace{1em}\noindent\textbf{Questão #1.}\hspace{0.4em}}{\par}
\newenvironment{solucao}{\vspace{0.3em}\noindent\textit{Solução.}\hspace{0.4em}\color{black!85}}{\par\vspace{0.5em}}


\usepackage[utf8]{inputenc}
\usepackage[T1]{fontenc}
\usepackage[brazil]{babel}
\usepackage{fullpage}
\usepackage{amsmath,amssymb}
\usepackage{enumitem}
\usepackage{xcolor}
\usepackage{fancyhdr}
\usepackage{titlesec}
\usepackage{listings}
\usepackage{hyperref}
\usepackage{booktabs}
\usepackage{graphicx}

\definecolor{matd48azul}{HTML}{0B4F6C}
\definecolor{matd48verde}{HTML}{2E8B57}

\titleformat{\section}{\normalfont\Large\bfseries\color{matd48azul}}{\thesection}{1em}{}
\titleformat{\subsection}{\normalfont\large\bfseries\color{matd48azul}}{\thesubsection}{1em}{}

\providecommand{\ListaTitulo}{Lista de Exercícios}

\setlength{\headheight}{14pt}
\pagestyle{fancy}
\fancyhf{}
\lhead{\small MATD48 --- Planejamento de Experimentos A}
\rhead{\small \ListaTitulo}
\cfoot{\thepage}
\renewcommand{\headrulewidth}{0.4pt}

\lstset{
  basicstyle=\ttfamily\small,
  breaklines=true,
  frame=single,
  columns=fullflexible,
  backgroundcolor=\color{gray!8},
  commentstyle=\color{matd48verde},
  keywordstyle=\color{matd48azul}\bfseries,
  language=R,
  extendedchars=true,
  literate=%
    {á}{{\'a}}1 {à}{{\`a}}1 {â}{{\^a}}1 {ã}{{\~a}}1
    {é}{{\'e}}1 {ê}{{\^e}}1 {í}{{\'i}}1 {ó}{{\'o}}1
    {ô}{{\^o}}1 {õ}{{\~o}}1 {ú}{{\'u}}1 {ç}{{\c c}}1
    {Á}{{\'A}}1 {À}{{\`A}}1 {Â}{{\^A}}1 {Ã}{{\~A}}1
    {É}{{\'E}}1 {Ê}{{\^E}}1 {Í}{{\'I}}1 {Ó}{{\'O}}1
    {Ô}{{\^O}}1 {Õ}{{\~O}}1 {Ú}{{\'U}}1 {Ç}{{\c C}}1
}

\hypersetup{colorlinks=true, linkcolor=matd48azul, urlcolor=matd48azul, citecolor=matd48azul}

\newcommand{\cabecalho}[2]{%
  \begin{center}
    {\Large\bfseries\color{matd48azul} #1}\\[0.3em]
    {\large #2}\\[0.3em]
    \rule{\linewidth}{0.6pt}
  \end{center}
  \vspace{0.5em}
}

\newenvironment{questao}[1]{\vspace{1em}\noindent\textbf{Questão #1.}\hspace{0.4em}}{\par}
\newenvironment{solucao}{\vspace{0.3em}\noindent\textit{Solução.}\hspace{0.4em}\color{black!85}}{\par\vspace{0.5em}}


\begin{document}

\cabecalho{Lista de Exercícios 14}{Fatoriais $3^k$, fracionados $2^{k-p}$ e superfície de resposta (Capítulo 6 / Aula 14)}

\noindent \textit{Entrega: até a próxima aula. Justifique todas as respostas --- nestas listas, a
justificativa vale mais que a resposta final. O gabarito completo está em \texttt{Gabarito14.pdf}.}

\begin{questao}{1} \textbf{(Engenharia --- estrutura de graus de liberdade de um $3^2$)}

Um experimento de usinagem cruza velocidade de corte (3 níveis) com ângulo de saída da ferramenta
(3 níveis) sobre a energia de corte consumida, com $r=4$ repetições por combinação
($n=3\times3\times4=36$).

\begin{enumerate}[label=(\alph*)]
  \item Quantos graus de liberdade tem cada efeito principal? E a interação $A{\times}B$?
    Justifique a partir da regra geral de um $3^k$ (Seção 6.4 do livro-texto).
  \item Monte a tabela de graus de liberdade completa (Velocidade, Ângulo, Velocidade$\times$
    Ângulo, Erro, Total) para este desenho específico ($a=b=3$, $r=4$).
  \item Um fatorial $2^2$ nos mesmos dois fatores (só 2 níveis cada) jamais conseguiria estimar um
    tipo específico de efeito que o $3^2$ estima. Qual é, e por quê?
\end{enumerate}
\end{questao}

\begin{questao}{2} \textbf{(Dedução --- estrutura de aliases de um fracionado $2^{4-1}$)}

Considere um fatorial fracionado $2^{4-1}$ com relação geradora $D=ABC$, equivalente à relação
definidora $I=ABCD$.

\begin{enumerate}[label=(\alph*)]
  \item Usando a propriedade $A^2=B^2=C^2=D^2=I$, calcule o alias de $A$ multiplicando ambos os
    lados de $I=ABCD$ por $A$: simplifique $A\cdot ABCD$ até obter uma expressão do tipo
    "$A$ é aliasado com \underline{\hspace{2cm}}".
  \item Repita o procedimento para encontrar o alias de $AB$ (multiplique $I=ABCD$ por $AB$).
  \item Repita para encontrar o alias de $C$.
  \item Classifique este fracionado quanto à \textbf{resolução} (III, IV ou V) e explique, em uma
    frase, o que essa resolução garante sobre a limpeza dos efeitos principais.
\end{enumerate}
\end{questao}

\begin{questao}{3} \textbf{(Ciência de dados/Engenharia --- ponto estacionário de uma superfície de resposta)}

Um modelo de segunda ordem foi ajustado para o rendimento de um processo, em função de duas
variáveis codificadas $x_1$ e $x_2$:
$$
\hat y = 80 + 4x_1 - 6x_2 - 2x_1^2 - 3x_2^2 + 2x_1x_2.
$$

\begin{enumerate}[label=(\alph*)]
  \item Calcule $\partial \hat y/\partial x_1$ e $\partial \hat y/\partial x_2$.
  \item Iguale as duas derivadas parciais a zero e resolva o sistema linear resultante para
    encontrar o ponto estacionário $(x_1^*, x_2^*)$.
  \item Escreva a matriz Hessiana (segundas derivadas) do modelo e calcule seus autovalores.
  \item Classifique o ponto estacionário (máximo, mínimo ou sela) a partir dos autovalores, e
    calcule $\hat y$ no ponto estacionário.
\end{enumerate}
\end{questao}

\begin{questao}{4} \textbf{(Confusão em $3^k$ --- comparando os geradores $AB$ e $AB^2$)}

Para um fatorial $3^2$ (fatores $A$ e $B$, níveis codificados $0,1,2$), o livro-texto usa o
gerador $AB^2$ (isto é, $L=A+2B \bmod 3$) para dividir as 9 combinações em 3 blocos de 3.

\begin{enumerate}[label=(\alph*)]
  \item Construa a tabela completa de $L_{AB^2} = (A+2B)\bmod3$ para as 9 combinações de $(A,B)$.
  \item Construa agora a tabela de $L_{AB} = (A+B)\bmod3$ para as mesmas 9 combinações.
  \item Compare as duas partições em blocos. A combinação $(A{=}0,B{=}0)$ e a combinação
    $(A{=}1,B{=}1)$ ficam no mesmo bloco sob $AB^2$? E sob $AB$? O que essa diferença mostra sobre
    $AB$ e $AB^2$ serem contrastes \emph{distintos} (não intercambiáveis)?
\end{enumerate}
\end{questao}

\begin{questao}{5} \textbf{(Ciência de dados/Engenharia --- exercício computacional em R)}

O código abaixo ajusta o modelo de segunda ordem da Questão 3-style aos dados reais de energia de
corte e varre uma grade para localizar o mínimo predito na região experimental. Está incompleto.

\begin{lstlisting}
energia <- read.csv("energia.csv")

# (a) complete a formula do modelo de segunda ordem
modelo <- lm(______________________, data = energia)

# (b) construa uma grade de 100 x 100 pontos dentro do intervalo observado
#     de Velocidad e angulo, e obtenha a energia predita em cada ponto
grade <- expand.grid(
  Velocidad = seq(min(energia$Velocidad), max(energia$Velocidad), length.out = 100),
  angulo    = seq(______________________, ______________________, length.out = 100)
)
grade$pred <- predict(______________________, newdata = grade)

# (c) identifique a linha da grade com a menor energia predita
grade[______________________, ]
\end{lstlisting}

\begin{enumerate}[label=(\alph*)]
  \item Complete as quatro linhas indicadas.
  \item O ponto estacionário do modelo ajustado a estes dados (encontrado resolvendo o gradiente
    igual a zero, como na Questão 3) cai fora da região onde a energia predita é mínima na grade.
    Que tipo de ponto estacionário (máximo, mínimo ou sela) é consistente com essa observação?
\end{enumerate}
\end{questao}

\end{document}
